\subsection{Метод условного градиента}
\label{subsec:cond-gradient}

Это тот же метод возможных направлений, но относящийся ещё и к методам
\emph{линейной аппроксимации}.

\[
  \begin{cases}
    \varphi(x) \to \min , & x \in S , \\
    S \ \text{— выпуклый компакт} ,
  \end{cases}
\]
$\varphi$ — выпуклая, непрерывно дифференцируемая, с липшицевым градиентом
(константа $R$).

\subsubsection*{Алгоритм}

Выбираются $\alpha_0 \in (0,1)$ (иногда $\alpha_0 = 1$), $\lambda \in (0,1)$,
$x_0 \in S$. Строится последовательность $\{x_k\} \subset S$.

Пусть выполнен $k$-й шаг и $x_k$ построена. Решается \textbf{вспомогательная
задача}
\[
  \nabla\T\varphi(x_k)\, x \ \longrightarrow\ \min_{x \in S} .
\]
Пусть $y_k$ — её решение. Выбираем
\[
  s_k = y_k - x_k \quad\text{— направление спуска},
  \qquad
  x_{k+1} = x_k + \alpha_k s_k = x_k + \alpha_k (y_k - x_k) ,
\]
и обозначаем
\[
  \eta_k = \nabla\T\varphi(x_k)\,(y_k - x_k) .
\]

\begin{idea}
Вспомогательная задача — это минимизация \emph{линейной} аппроксимации
$\varphi$ (постоянное слагаемое $\varphi(x_k) - \nabla\T\varphi(x_k)x_k$ на
точку минимума не влияет). Отсюда два практических замечания, помеченные на
лекции стрелками:
\begin{itemize}[nosep,left=0pt]
  \item если $S$ задано линейными ограничениями, вспомогательная задача — это
        задача ЛП, и решается симплекс-методом (\S4);
  \item в любом случае минимум линейной функции на выпуклом компакте
        достигается на границе (и, если $S$ — многогранник, в вершине).
        Поэтому $y_k$ — всегда «крайняя» точка, и шаг $\alpha_k$ обязателен:
        идти в $y_k$ целиком, как правило, нельзя.
\end{itemize}
\end{idea}

\inputfig{figures/fig-frank-wolfe}

Так как $y_k$ минимизирует $\nabla\T\varphi(x_k)x$ по $S$, а $x_k \in S$,
\[
  \nabla\T\varphi(x_k)\, y_k \ \le\ \nabla\T\varphi(x_k)\, x_k
  \qquad\Longrightarrow\qquad
  \eta_k \le 0 .
\]
Возможны два случая:
\begin{enumerate}[label=\arabic*),left=0pt]
  \item $\eta_k < 0$ \ $\Longrightarrow$ \
        $\eta_k = \nabla\T\varphi(x_k)\, s_k < 0$, то есть $s_k$ — направление
        спуска;
  \item $\eta_k = 0$ \ $\Longrightarrow$ \ движение прекращается (см. условие
        останова ниже).
\end{enumerate}
Относительно $x_{k+1}$ требуется: 2.1)~$x_{k+1}$ допустима; 2.2)~$\{x_k\}$ —
последовательность спуска. Обоим требованиям и подчинён выбор $\alpha_k$.

\subsubsection*{Допустимость}

Пусть $\eta_k < 0$. Тогда
\[
  x_{k+1} = x_k + \alpha_k (y_k - x_k) = (1 - \alpha_k) x_k + \alpha_k y_k ,
\]
и так как $x_k \in S$, $y_k \in S$, при $\alpha_k \in [0,1]$ точка $x_{k+1}$
лежит в $S$ — множество $S$ выпукло.

\subsubsection*{Выбор шага}

По формуле конечных приращений
\[
  \varphi(x_k + \alpha_k s_k)
  = \varphi(x_k) + \alpha_k \nabla\T\varphi(\bar{x}_k)\, s_k ,
  \qquad
  \bar{x}_k = x_k + \theta \alpha_k s_k , \quad \theta \in [0,1] .
\]
Прибавляя и вычитая $\alpha_k \nabla\T\varphi(x_k) s_k$ и пользуясь
липшицевостью градиента,
\begin{align*}
  \varphi(x_k + \alpha_k s_k)
  &= \varphi(x_k) + \alpha_k \underbrace{\nabla\T\varphi(x_k) s_k}_{=\ \eta_k}
     + \alpha_k \big(\nabla\varphi(\bar{x}_k) - \nabla\varphi(x_k)\big)\T s_k \\
  &\le \varphi(x_k) + \alpha_k \eta_k
     + \alpha_k \norm{\nabla\varphi(\bar{x}_k) - \nabla\varphi(x_k)}
       \cdot \norm{s_k} \\
  &\le \varphi(x_k) + \alpha_k \eta_k
     + \alpha_k R \underbrace{\norm{\bar{x}_k - x_k}}_{=\ \theta\alpha_k\norm{s_k}}
       \cdot \norm{s_k}
   \ \le\ \varphi(x_k) + \alpha_k \eta_k + \alpha_k^2 R \norm{s_k}^2 .
\end{align*}
Потребуем
\begin{equation}\label{eq:fw-req}
  \varphi(x_k) + \alpha_k \eta_k + \alpha_k^2 R \norm{s_k}^2
  \ \le\ \varphi(x_k) + \tfrac12 \alpha_k \eta_k ,
\end{equation}
то есть $\tfrac12 \alpha_k \eta_k + \alpha_k^2 R \norm{s_k}^2 \le 0$, откуда
\begin{equation}\label{eq:fw-alpha}
  \alpha_k \ \le\ -\frac{1}{2}\, \eta_k \cdot \frac{1}{R \norm{s_k}^2} .
\end{equation}

\begin{remark}
Пусть $x_0$ допустима, построена $\{x_k\} \subset S$, и правило выбора шага
привело к~\eqref{eq:fw-alpha}. Тогда
\begin{equation}\label{eq:fw-1}
  \varphi(x_{k+1}) \ \le\ \varphi(x_k) + \tfrac12 \alpha_k \eta_k .
\end{equation}
\end{remark}

\begin{idea}
При $\eta_k = 0$ правая часть~\eqref{eq:fw-alpha} равна нулю: шага не делается,
и это корректно — именно $\eta_k = 0$ служит признаком остановки.

При $\eta_k < 0$ имеем $s_k \ne 0$, и шаг $\alpha_k$ может быть выбран.
Берём его по принципу дробления, $\alpha_k = \alpha_0 \lambda^{i_k}$: это первое
значение, удовлетворяющее~\eqref{eq:fw-alpha}. Тогда предыдущее пробное
значение $\alpha = \alpha_k / \lambda$ условию~\eqref{eq:fw-alpha} уже
\emph{не} удовлетворяло (само неравенство~\eqref{eq:fw-req} от выбора $\alpha$
не зависит по форме), то есть
\[
  \alpha_0 \lambda^{\,i_k - 1} \ >\ -\frac{1}{2}\,\eta_k \cdot
  \frac{1}{R\norm{s_k}^2}
  \qquad\Longrightarrow\qquad
  \alpha_k \ >\ -\frac{1}{2}\,\eta_k \cdot \frac{\lambda}{R\norm{s_k}^2} .
\]
Наконец, $\norm{s_k} = \norm{y_k - x_k}$, где $y_k, x_k \in S$, а $S$ —
компакт, поэтому существует $C$: $\norm{s_k} < C$. Итого
\begin{equation}\label{eq:fw-2}
  \alpha_k \ >\ -\frac{\lambda}{2RC^2}\, \eta_k .
\end{equation}
\end{idea}

\subsubsection*{Основная лемма}

\begin{lemma}\label{lem:fw}
Пусть $\{x_k\}$ построена по методу условного градиента. Тогда
\[
  1^\circ\ \{x_k\} \subset S ;
  \qquad
  2^\circ\ \varphi(x_k) \ \text{монотонно убывает} ;
  \qquad
  3^\circ\ \eta_k \xrightarrow[k \to \infty]{} 0 .
\]
\end{lemma}

\begin{proof}
$1^\circ$ доказано выше (выпуклость $S$ и $\alpha_k \in [0,1]$).

\medskip
$2^\circ$--$3^\circ$. Из~\eqref{eq:fw-1} и~\eqref{eq:fw-2}, учитывая
$\eta_k < 0$ (умножение на отрицательное число меняет знак неравенства),
\begin{equation}\label{eq:fw-3}
  \varphi(x_{k+1}) - \varphi(x_k)
  = \varphi(x_k + \alpha_k s_k) - \varphi(x_k)
  \ \le\ \tfrac12 \alpha_k \eta_k
  \ <\ -\frac{\lambda}{4RC^2}\, \eta_k^2 \ \le\ 0 ,
\end{equation}
то есть $\varphi(x_k)$ убывает. Просуммируем~\eqref{eq:fw-3} по
$k = \overline{0, m-1}$; слагаемые телескопируются:
\[
  \varphi(x_m) - \varphi(x_0)
  \ <\ -\frac{\lambda}{4RC^2} \sum_{k=0}^{m-1} \eta_k^2 .
\]
Так как $\varphi(x_m) \ge \varphi(x_*)$,
\[
  \varphi(x_*) - \varphi(x_0) \ \le\ \varphi(x_m) - \varphi(x_0)
  \ <\ -\frac{\lambda}{4RC^2} \sum_{k=0}^{m-1} \eta_k^2 ,
\]
откуда при $m \to +\infty$
\[
  \sum_{k=0}^{\infty} \eta_k^2
  \ \le\ \frac{4RC^2}{\lambda}\big(\varphi(x_0) - \varphi(x_*)\big) < +\infty .
\]
Ряд из $\eta_k^2$ сходится, значит $\eta_k^2 \to 0$ и $\eta_k \to 0$.
\end{proof}

\subsubsection*{Условие останова}

\begin{enumerate}[label=\arabic*),left=0pt]
  \item за конечное число шагов получено $\eta_k = 0$;
  \item последовательность $\{\varphi(x_k)\}$ перестала заметно убывать.
\end{enumerate}

Разберём первый случай. Пусть $\eta_k = 0$, то есть
$\nabla\T\varphi(x_k)(y_k - x_k) = 0$, значит
$\nabla\T\varphi(x_k) y_k = \nabla\T\varphi(x_k) x_k$. Но $y_k$ по построению
минимизирует линейную функцию:
\[
  \nabla\T\varphi(x_k)\, y_k \ \le\ \nabla\T\varphi(x_k)\, x
  \qquad \forall x \in S ,
\]
следовательно $\nabla\T\varphi(x_k) x \ge \nabla\T\varphi(x_k) x_k$, то есть
\[
  \nabla\T\varphi(x_k)\,(x - x_k) \ \ge\ 0 \qquad \forall x \in S .
\]
Это условие экстремума выпуклой функции на выпуклом множестве — значит $x_k$
оптимальна.

\subsubsection*{Теорема о сходимости}

\begin{theorem}[о сходимости метода условного градиента]\label{thm:fw}
Пусть выполнены условия применимости метода и $\{x_k\}$ построена по методу
условного градиента. Тогда в любой предельной точке последовательности
выполняется условие минимума задачи $\varphi(x) \to \min$, $x \in S$.
\end{theorem}

\begin{proof}
Пусть $\bar{x}_*$ — предельная точка $\{x_k\}$, то есть существует
подпоследовательность $x_{k_j} \to \bar{x}_*$. По лемме~\ref{lem:fw}
\[
  \eta_{k_j} = \nabla\T\varphi(x_{k_j})\,(y_{k_j} - x_{k_j})
  \xrightarrow[j \to \infty]{} 0 .
\]
Последовательность $\{y_{k_j}\}$ лежит в компакте $S$, поэтому, переходя при
необходимости к подпоследовательности, считаем $y_{k_j} \to y_*$. По
непрерывности $\nabla\varphi$
\begin{equation}\label{eq:fw-4}
  \nabla\T\varphi(\bar{x}_*)\,(y_* - \bar{x}_*) = 0 .
\end{equation}

Далее — так же, как при разборе условия останова. Для всех $x \in S$
\[
  \nabla\T\varphi(x_{k_j})\, y_{k_j} \ \le\ \nabla\T\varphi(x_{k_j})\, x ,
\]
и при $k_j \to \infty$
\[
  \nabla\T\varphi(\bar{x}_*)\, y_* \ \le\ \nabla\T\varphi(\bar{x}_*)\, x .
\]
По~\eqref{eq:fw-4} левая часть равна
$\nabla\T\varphi(\bar{x}_*)\,\bar{x}_*$, поэтому
\[
  \nabla\T\varphi(\bar{x}_*)\, \bar{x}_* \ \le\ \nabla\T\varphi(\bar{x}_*)\, x
  \qquad\Longrightarrow\qquad
  \nabla\T\varphi(\bar{x}_*)\,(x - \bar{x}_*) \ \ge\ 0 \quad \forall x \in S .
\]
Это условие экстремума в точке $\bar{x}_*$; в силу выпуклости $\varphi$ оно
достаточно, значит $\bar{x}_*$ — точка минимума. Отсюда
\[
  \lim_{k \to \infty} \varphi(x_k) = \varphi(x_*) ,
  \qquad \varphi(x_*) = \min_{x \in S} \varphi(x) .
  \qedhere
\]
\end{proof}

\begin{remark}
Кроме того, справедлива оценка скорости
\[
  \varphi(x_k) - \varphi(x_*) \ \le\ \frac{D}{k} ,
  \qquad D > 0 \ \text{— const} .
\]
\end{remark}

\begin{idea}
Оценка $O(1/k)$ — классический результат для метода условного градиента
(Франка--Вульфа); на лекции она приведена без доказательства. Константа $D$
выражается через диаметр $S$ и константу Липшица градиента:
$D \sim R\,\operatorname{diam}^2 S$. Важно, что скорость здесь \emph{сублинейная}
— заметно хуже геометрической сходимости методов из главы о безусловной
минимизации. Платой за это является крайняя дешевизна итерации: вместо
проектирования на $S$ (как в методе Розена) решается одна линейная задача.
\end{idea}
